% Slide 23 — Fechamento
\begin{frame}{Fechamento}
  \begin{center}
    \large
    \textcolor{darkblue}{\textbf{Seleção contextual mostra ganho frente à política única no recorte avaliado}}\\[0.35em]
    \normalsize A decisão de reposição melhora quando considera o contexto
    operacional de cada série.
  \end{center}
  \vspace{0.75em}
  \begin{columns}[T,onlytextwidth]
    \column{0.33\textwidth}
      \centering\small\textbf{Seleção contextual}\\
      \footnotesize evidência empírica + PSE como meta-modelo
    \column{0.33\textwidth}
      \centering\small\textbf{Pipeline reprodutível}\\
      \footnotesize Kedro, rastreabilidade completa
    \column{0.33\textwidth}
      \centering\small\textbf{Dados reais brasileiros}\\
      \footnotesize PB e BA, regime \textit{Lumpy}
  \end{columns}
  \vspace{0.75em}
  \begin{center}
    {\color{darkblue}\rule{0.55\linewidth}{0.9pt}}
  \end{center}
  \vspace{0.4em}
  \begin{center}
    \LARGE\textcolor{darkblue}{\textbf{Obrigado.}}\\[0.25em]
    \normalsize\textcolor{slategray}{Perguntas e discussão}
  \end{center}
  \note{
    Para concluir: a contribuição desta dissertação não é um novo algoritmo
    de otimização, mas a formulação e avaliação empírica de um framework de
    seleção contextual de políticas de reposição -- o AIPE -- apoiado em um
    pipeline reprodutível e em dados reais de uma rede varejista brasileira.
    Os dois experimentos executados, na Paraíba e na Bahia, mostram
    heterogeneidade genuína entre políticas e uma primeira evidência de que
    selecionar por perfil produz diferença econômica mensurável, com
    trade-offs que precisam ser geridos explicitamente. A expansão nacional
    e o treinamento do PSE em escala são os próximos passos. Fico à
    disposição para perguntas.
  }
\end{frame}
